\chapter{2020年全国III卷(文)}

\section{单选题}

\begin{problem}
已知集合 \(A=\{1,2,3,5,7,11\}\),\(B=\{x\mid 3<x<15\}\),则 \(A\cap B\) 中元素的个数为(\quad)

\choices
    {\(2\)}
    {\(3\)}
    {\(4\)}
    {\(5\)}
\end{problem}

\begin{problem}
若 \(\overline{z}(1+\i)=1-\i\),则 \(z=\)(\quad)

\choices
    {\(1-\i\)}
    {\(1+\i\)}
    {\(-\i\)}
    {\(\i\)}
\end{problem}

\begin{problem}
设一组样本数据 \(x_1\),\(x_2\),\(\cdots\),\(x_n\) 的方差为 \(0.01\),则数据 \(10x_1\),\(10x_2\),\(\cdots\),\(10x_n\) 的方差为(\quad)

\choices
    {\(0.01\)}
    {\(0.1\)}
    {\(1\)}
    {\(10\)}
\end{problem}

\begin{problem}
Logistic 模型是常用数学模型之一,可应用于流行病学领域. 有学者根据公布数据建立了某地区新冠肺炎累计确诊病例数 \(I(t)\)(\(t\) 的单位:天)的 Logistic 模型:\(I(t)=\dfrac{K}{1+\e^{-0.23(t-53)}}\),其中 \(K\) 为最大确诊病例数. 当 \(I(t^*)=0.95K\) 时,标志着已初步遏制疫情,则 \(t^*\) 约为(\quad)(\(\ln 19\approx 3\))

\choices
    {\(60\)}
    {\(63\)}
    {\(66\)}
    {\(69\)}
\end{problem}

\begin{problem}
已知 \(\sin\theta+\sin\left( \theta+\dfrac{\pi}{3} \right)=1\),则 \(\sin\left( \theta+\dfrac{\pi}{6} \right)=\)(\quad)

\choices
    {\(\dfrac{1}{2}\)}
    {\(\dfrac{\sqrt3}{3}\)}
    {\(\dfrac{2}{3}\)}
    {\(\dfrac{\sqrt2}{2}\)}
\end{problem}

\begin{problem}
在平面内,\(A\),\(B\) 是两个定点,\(C\) 是动点. 若 \(\overrightarrow{AC}\cdot\overrightarrow{BC}=1\),则点 \(C\) 的轨迹为(\quad)

\choices
    {圆}
    {椭圆}
    {抛物线}
    {直线}
\end{problem}

\begin{problem}
设 \(O\) 为坐标原点,直线 \(x=2\) 与抛物线 \(C:y^2=2px\)(\(p>0\))交于 \(D\),\(E\) 两点,若 \(OD\perp OE\),则 \(C\) 的焦点坐标为(\quad)

\choices
    {\(\left( \dfrac{1}{4},0 \right)\)}
    {\(\left( \dfrac{1}{2},0 \right)\)}
    {\((1,0)\)}
    {\((2,0)\)}
\end{problem}

\begin{problem}
点 \((0,-1)\) 到直线 \(y=k(x+1)\) 距离的最大值为(\quad)

\choices
    {\(1\)}
    {\(\sqrt2\)}
    {\(\sqrt3\)}
    {\(2\)}
\end{problem}

\begin{problem}
如图为某几何体的三视图,则该几何体的表面积是(\quad)

\bitmapfigure[width=5.4cm]{img/2020/national_paper_3_liberal/q9_fig1.png}

\choices
    {\(6+4\sqrt2\)}
    {\(4+4\sqrt2\)}
    {\(6+2\sqrt3\)}
    {\(4+2\sqrt3\)}
\end{problem}

\begin{problem}
设 \(a=\log_3 2\),\(b=\log_5 3\),\(c=\dfrac{2}{3}\),则(\quad)

\choices
    {\(a<c<b\)}
    {\(a<b<c\)}
    {\(b<c<a\)}
    {\(c<a<b\)}
\end{problem}

\begin{problem}
在 \(\triangle ABC\) 中,\(\cos C=\dfrac{2}{3}\),\(AC=4\),\(BC=3\),则 \(\tan B=\)(\quad)

\choices
    {\(\sqrt5\)}
    {\(2\sqrt5\)}
    {\(4\sqrt5\)}
    {\(8\sqrt5\)}
\end{problem}

\begin{problem}
已知函数 \(f(x)=\sin x+\dfrac{1}{\sin x}\),则(\quad)

\choices
    {\(f(x)\) 的最小值为 \(2\)}
    {\(f(x)\) 的图象关于 \(y\) 轴对称}
    {\(f(x)\) 的图象关于直线 \(x=\pi\) 对称}
    {\(f(x)\) 的图象关于直线 \(x=\dfrac{\pi}{2}\) 对称}
\end{problem}

\section{填空题}

\begin{problem}
若 \(x\),\(y\) 满足约束条件
\(\begin{cases}
x+y\ge 0,\\
2x-y\ge 0,\\
x\le 1
\end{cases}\)
则 \(z=3x+2y\) 的最大值为\fillinblank{}.
\end{problem}

\begin{problem}
设双曲线 \(C:\dfrac{x^2}{a^2}-\dfrac{y^2}{b^2}=1\)(\(a>0\),\(b>0\))的一条渐近线为 \(y=\sqrt2x\),则 \(C\) 的离心率为\fillinblank{}.
\end{problem}

\begin{problem}
设函数 \(f(x)=\dfrac{\e^x}{x+a}\),若 \(f'(1)=\dfrac{\e}{4}\),则 \(a=\fillinblank{}\).
\end{problem}

\begin{problem}
已知圆锥的底面半径为 \(1\),母线长为 \(3\),则该圆锥内半径最大的球的体积为\fillinblank{}.
\end{problem}

\section{解答题}

\begin{problem}
设等比数列 \(\{a_n\}\) 满足 \(a_1+a_2=4\),\(a_3-a_1=8\).

\begin{enumerate}
    \item 求 \(\{a_n\}\) 的通项公式;
    \item 记 \(S_n\) 为数列 \(\{\log_3 a_n\}\) 的前 \(n\) 项和. 若 \(S_m+S_{m+1}=S_{m+3}\),求 \(m\).
\end{enumerate}
\end{problem}

\begin{problem}
某学生兴趣小组随机调查了某市 \(100\) 天中每天的空气质量等级和当天到某公园锻炼的人次,整理数据得到下表(单位:天):

\FigureLayoutDeclare{img/2020/national_paper_3_liberal/q18_fig1.png}{6.5cm}{4bp 0bp 0bp 15bp}

\newcommand{\airqualityheader}{%
    \choicebitmap[width=6.5cm,trim=4bp 0bp 0bp 15bp,clip]{img/2020/national_paper_3_liberal/q18_fig1.png}%
}

\begin{center}
\begin{tabular}{c|ccc}
\airqualityheader & \([0,200]\) & \((200,400]\) & \((400,600]\) \\
\hline
\(1\)(优) & \(2\) & \(16\) & \(25\) \\
\(2\)(良) & \(5\) & \(10\) & \(12\) \\
\(3\)(轻度污染) & \(6\) & \(7\) & \(8\) \\
\(4\)(中度污染) & \(7\) & \(2\) & \(0\)
\end{tabular}
\end{center}

\begin{enumerate}
    \item 分别估计该市一天的空气质量等级为 \(1\),\(2\),\(3\),\(4\) 的概率;
    \item 求一天中到该公园锻炼的平均人次的估计值(同一组中的数据用该组区间的中点值为代表);
    \item 若某天的空气质量等级为 \(1\) 或 \(2\),则称这天``空气质量好'';若某天的空气质量等级为 \(3\) 或 \(4\),则称这天``空气质量不好''. 根据所给数据,完成下面的 \(2\times 2\) 列联表,并根据列联表,判断是否有 \(95\%\) 的把握认为一天中到该公园锻炼的人次与该市当天的空气质量有关?
\end{enumerate}

\begin{center}
\begin{tabular}{c|cc}
 & 人次\(\le 400\) & 人次\(>400\) \\
\hline
空气质量好 &  &  \\
空气质量不好 &  &
\end{tabular}
\end{center}

附:\(K^2=\dfrac{n(ad-bc)^2}{(a+b)(c+d)(a+c)(b+d)}\).

\begin{center}
\begin{tabular}{c|ccc}
\(P(K^2\ge k)\) & \(0.050\) & \(0.010\) & \(0.001\) \\
\hline
\(k\) & \(3.841\) & \(6.635\) & \(10.828\)
\end{tabular}
\end{center}
\end{problem}

\begin{problem}
如图,在长方体 \(ABCD-A_1B_1C_1D_1\) 中,点 \(E\),\(F\) 分别在棱 \(DD_1\),\(BB_1\) 上,且 \(2DE=ED_1\),\(BF=2FB_1\). 证明:

\begin{enumerate}
    \item 当 \(AB=BC\) 时,\(EF\perp AC\);
    \item 点 \(C_1\) 在平面 \(AEF\) 内.
\end{enumerate}

\bitmapfigure[width=6.0cm]{img/2020/national_paper_3_liberal/q19_fig1.png}
\end{problem}

\begin{problem}
已知函数 \(f(x)=x^3-kx+k^2\).

\begin{enumerate}
    \item 讨论 \(f(x)\) 的单调性;
    \item 若 \(f(x)\) 有三个零点,求 \(k\) 的取值范围.
\end{enumerate}
\end{problem}

\begin{problem}
已知椭圆 \(C:\dfrac{x^2}{25}+\dfrac{y^2}{m^2}=1\)(\(0<m<5\))的离心率为 \(\dfrac{\sqrt{15}}{4}\),\(A\),\(B\) 分别为 \(C\) 的左,右顶点.

\begin{enumerate}
    \item 求 \(C\) 的方程;
    \item 若点 \(P\) 在 \(C\) 上,点 \(Q\) 在直线 \(x=6\) 上,且 \(|BP|=|BQ|\),\(BP\perp BQ\),求 \(\triangle APQ\) 的面积.
\end{enumerate}
\end{problem}

\section{选考题:解答题}

\begin{problem}
在直角坐标系 \(xOy\) 中,曲线 \(C\) 的参数方程为
\(\begin{cases}
x=2-t-t^2,\\
y=2-3t+t^2
\end{cases}\)
(\(t\) 为参数且 \(t\neq 1\)),\(C\) 与坐标轴交于 \(A\),\(B\) 两点.

\begin{enumerate}
    \item 求 \(|AB|\);
    \item 以坐标原点为极点,\(x\) 轴正半轴为极轴建立极坐标系,求直线 \(AB\) 的极坐标方程.
\end{enumerate}
\end{problem}

\begin{problem}
设 \(a\),\(b\),\(c\in \R\),\(a+b+c=0\),\(abc=1\).

\begin{enumerate}
    \item 证明:\(ab+bc+ca<0\);
    \item 用 \(\max\{a,b,c\}\) 表示 \(a\),\(b\),\(c\) 中的最大值,证明:\(\max\{a,b,c\}\ge \sqrt[3]{4}\).
\end{enumerate}
\end{problem}

